\chapter{Research Method}

\section{Artifact Types}

The project uses three durable research artifact types, defined in
\path{docs/research/000-research-index.md}.

\begin{description}
  \item[RL] Research logs. These are living synthesis documents organized by
  question and decision impact.
  \item[SR] Source reviews. These summarize external public documentation,
  open-source implementation behavior, reverse-engineering notes, and support
  limits.
  \item[EX] Experiments. These are controlled probes with scripts, generated
  artifacts, expected outcomes, observed outcomes, and a distilled conclusion.
\end{description}

Raw command output belongs in experiment artifacts. Durable conclusions belong
in the experiment README and then in the affected research logs.

\section{Staged Gates}

The staged gates are:

\begin{description}
  \item[Gate A] Minimally correct v1 parser: checkpoint selection, object
  resolution, FS tree records, namespace, logical size, and validation.
  \item[Gate B] Support boundary: runtime source classes, encryption, feature
  gates, and fallback.
  \item[Gate C] Safe incremental scanning: identity, subtree reuse, persistence,
  and invalidation.
  \item[Gate D] Product semantics and optimization: snapshots, firmlinks,
  volume groups, and performance model.
\end{description}

This ordering prevents performance work from outrunning correctness and prevents
product claims from outrunning support evidence.

\section{Claim Discipline}

Every durable claim should be tagged mentally as one of:

\begin{itemize}
  \item \textbf{Spec}: backed by Apple/public documentation.
  \item \textbf{Observation}: backed by repo experiments or converging external
  implementations.
  \item \textbf{Hypothesis}: plausible, useful, but not yet proven.
\end{itemize}

The manual should preserve that discipline. It is acceptable for a chapter to
say "not yet known"; that is better than encoding a guess as architecture.
